% Online Appendix to Section 3 (sec:data). Specification: scripts/utils/raw_inputs.R
% (resolve_versioned_raw_input), config/raw_inputs_manifest.yml, data/raw/raw_inputs_resolved.yml,
% R/utils/country.R, R/utils/eia_intl.R, scripts/98_pull_comtrade_energy_dyads.R,
% scripts/40_build_index_vintages.R, docs/methodology.md. Final rerun of 12 October 2026.
\section{Data: releases, harmonisation and years}\label{oa:data}

\paragraph{The release rule.} Inputs that are published in successive releases are staged under
a file-name pattern, and the run selects one release by rule rather than by hand. Let $F_P$ be the
staged files matching the pattern of input $P$, $k(f)$ the vintage key parsed from a file's name
(a date or a year), $m(f)$ its modification time and $R(f)$ the number of distinct reporters in a
tabular extract. The release used is
\begin{equation}\label{eq:data_resolve}
  f^{*}_{P}=\operatorname*{lex\,max}_{f\in F_P:\ R(f)\ge\underline R_P}\ \bigl(k(f),\,m(f)\bigr),
\end{equation}
the newest declared vintage, with the modification time breaking ties and ordering unversioned
names, among the files that meet the coverage floor $\underline R_P$. The floor is 50 reporters for
the Comtrade critical-minerals extracts and absent elsewhere. Comtrade publishes a year's data as
reporters file, so the newest file can be a stub: the 2025 critical-minerals extracts have three
reporters each and are skipped in favour of the 2024 extracts (130 import, 128 export and 130
total-export reporters). Every run writes the releases it resolved, with their reporter counts, to
a vintage stamp (\texttt{data/raw/raw\_inputs\_resolved.yml}); Appendix~\ref{app:sources} reports
them. Inputs that are not pattern-resolved (the Comtrade dyads, the Red Book and USGS tables, the
EIA bulk file, CM-FDI, the ideal points and the project-authored lists) are single staged files
whose access dates are recorded in the input manifest.

\paragraph{Harmonisation of economies.} Every source is put on one canonical name per economy
before anything is ranked: the World Development Indicators name of the economy's ISO3 code,
matched through the \texttt{countrycode} package and a curated table of spellings. Taiwan, which the
World Bank's country table does not list, is added under its own code, including to the rosters of
the trade themes; Kosovo is carried as XKX. World, regional, income-group and residual rows
(``Other Asia Pacific'', ``Rest of world'') and historical entities are dropped from every theme
table, so that no aggregate is imputed or averaged as if it were an economy; theme builders also
remove aggregates from their comparison sets before ranking, while keeping them in world
denominators where the source defines the world total. Keys that two spellings of one economy
share are collapsed, with a warning if both carry different non-missing values. Sovereignty, which
the rankability rule uses, follows the Correlates of War state-system list \citep{cow2024state}:
195 states in 2024, the 193 UN members plus Kosovo and Taiwan.

\paragraph{Fossil production from EIA.} The Statistical Review lists fewer producers than consumers
and counts the rest in regional ``Other'' residuals. Where it has no production row for an economy
and fuel, the value is taken from the EIA International Energy Statistics \citep{eia2026international}
(bulk file of 3 October 2026) in the Review's units: oil as crude oil including lease condensate
plus natural-gas plant liquids, in thousand barrels a day; gas as dry gas in billion cubic metres
times 0.036 EJ; coal as terajoules (gross calorific value) times $0.95\times10^{-6}$ EJ. On economies
both sources report, the median ratio of the Review's value to the converted EIA value is 1.006 to
1.011 for oil, 0.97 to 0.99 for gas and 0.99 to 1.05 for coal. An EIA value is not used when EIA's
consumption of the fuel covers less than a tenth of the Review's, since the Review then counts under
the fuel something EIA excludes; in 2015--2025 this affects only Estonian coal (oil shale), which
stays missing. Where EIA's latest production year precedes the Review's consumption year (2025 gas
and coal from EIA 2024), self-sufficiency is computed with consumption of the same year and carried
forward. EIA values are observations, not imputations.

\paragraph{Reserves and uranium.} The USGS Mineral Commodity Summaries 2026 data release
\citep{usgs2026mcs} gives country reserves at end-2025; the report's version 1.3 erratum puts
Brazil's rare-earth reserves at 11~Mt (the Statistical Review's mineral sheet, an older USGS copy,
gives 21~Mt). A reported zero scores 0; a holder listed without a figure (NA, W or s) is missing in
its mineral variable, scored 0 and counted as imputed, and counts as zero in the technology
roll-ups, so world shares are measured on quantified holdings. ``Other countries'' and world-total
rows are never holders. The Red Book \citep{oecdnea2026uranium} supplies identified recoverable
resources (reasonably assured plus inferred) below USD 260/kgU at 1 January 2025 (Table 1.2a of
\emph{Uranium 2026}) and mine production (Table 1.18; \emph{Uranium 2026} for 2022--2024, earlier
editions for 2019--2021), completed with observed zeros over the economies of the Statistical
Review and the World Development Indicators.

\paragraph{Bilateral trade.} The all-reporter dyad file was pulled on 7--8 October 2026 for the
pre-review product list, with supplements on 9 October (the copper and aluminium inputs of the grid)
and 12 October (the new upstream codes). Comtrade's partner S19 (``Other Asia, n.e.s.'') is mapped
to Taiwan. Liberia, Mali and S\~ao Tom\'e and Pr\'incipe filed 2024 data after the pull and enter
through their partners' declarations. The country-level energy-trade file is derived from the
pull's World rows; an earlier direct extract, with 129--130 reporters, lacked 60 basket codes and
Qatar's exports of liquefied natural gas. Critical-minerals exports of non-reporting economies are
measured from partners' imports, moved from CIF to FOB with an HS-chapter ratio estimated on
importers that file both values.

\paragraph{Regional and staged inputs.} The IEA's 2035 announced manufacturing-capacity shares and
its electric-vehicle production and sales are regional for Europe, so EU members carry the EU's
value; economies outside the IEA's named regions take its ``Rest of world'' values. ILOSTAT
earnings are read from a copy staged on 30 September 2026, since the ILO's interface returns empty
responses to scripted requests. In the IMF lending-rate series a zero reading in a lending-type
rate is a placeholder and is dropped with the averages that include it: New Zealand's lending series
carries only zeros after January 2017, so its last genuine reading (4.76 percent) is used. Commodity
prices are in the units of the IMF series metadata.

\begin{table}[tbp]
\centering
\begin{threeparttable}
\caption{Years of the inputs in the final run}\label{tab:oa_data_years}
\footnotesize
\begin{tabularx}{\textwidth}{@{}>{\raggedright\arraybackslash}p{5.2cm} >{\raggedright\arraybackslash}X@{}}
\toprule
Source & Years used\\
\midrule
EI Statistical Review 2026 & Consumption and production to 2025 (growth 2020--2025); fossil
  reserves end-2020\\
EIA International Energy Statistics & Oil production to 2025; gas and coal to 2024\\
USGS Mineral Commodity Summaries 2026 & Reserves end-2025\\
Red Book (Uranium 2026) & Resources at 1 January 2025; production 2019--2024\\
IEA Critical Minerals Dataset 2026 & Base year and 2035 (Stated Policies)\\
IEA ETP 2024 capacity; IEA EV regional data & Announced capacity to 2035; EV production and sales
  2024\\
UN Comtrade dyads; energy-trade file & 2020 and 2024; 2024\\
Comtrade critical-minerals extracts & 2024 (2025 below the floor)\\
Atlas of Economic Complexity (HS92) & 2023\\
IEA WEO 2025; Global EV Outlook 2026 & 2024 and 2035; 2022, 2025 and 2035\\
BNEF New Energy Outlook 2024; Supply Chains 2025 & 2024 and 2035; 2025\\
IEA Clean Energy Technology Guide 2026 & Readiness 2020--2025\\
BNEF LCOE; IEA relative costs & 2025 and 2050; 2023\\
ILOSTAT; IMF lending rates, PPI, PCPS & Latest within ten years; to October, September and
  December 2025\\
Clean Investment Monitor & Investment 2018--2025; capacity 2025\\
IEA PAMS; Climate Action Tracker & Export of 22 January 2026; ratings of May 2022--June 2025\\
IMF Direct Investment Positions; CM-FDI & Mean of 2020--2024; edition of 5 October 2026\\
WDI; Doing Business; OECD CRS & 2022--2024 (GDP to 2025); 2019; 2013--2024\\
UNGA ideal points (Dataverse V39) & Sessions 74--78 (2019--2023)\\
\bottomrule
\end{tabularx}
\begin{tablenotes}[flushleft]\footnotesize
\item \textit{Notes:} Each component enters at the latest year its source reports
(equation~\eqref{eq:latest_year}). Appendix~\ref{app:sources} gives the releases.
\end{tablenotes}
\end{threeparttable}
\end{table}

\paragraph{Index vintages.} The headline indices are a single snapshot. A separate build re-cuts
each raw input to a requested year and reruns the same theme and pillar builders, so that the
pillars can be compared across years (Section~\ref{sec:robustness}). Only sources with a genuine
time dimension move: the Statistical Review, the Atlas trade data, IMF prices and rates, ILOSTAT
earnings and Clean Investment Monitor investment. Resource potential, reserves, the IEA critical
minerals release, BloombergNEF demand and costs, overcapacity and technology readiness are
snapshots or projections and are held fixed; a 2025 vintage uses Atlas 2023. The build records,
for each vintage, the share of each pillar's scored weight that re-derives and the year to which
each source resolved.
